\section{Proposed Method}
\label{sec:method}

\subsection{Pipeline Overview}

The pipeline is designed as a sequence of independent stages that exchange
data through JSON artifacts with defined schemas. This organization allows
each stage to be inspected, an individual stage to be rerun without
executing the entire pipeline, and different verification methods to be
compared on the same retrieval results.

The overall processing flow comprises five stages:

\begin{enumerate}
    \item normalizing legal text into condition--effect rules;

    \item constructing the legal causal graph and causal memory;

    \item retrieving events, rules, and multi-hop causal paths;

    \item analyzing claims from the query and performing verification with
    LegalSCM or path ablation;

    \item selecting evidence, determining the answer intent, and generating
    an answer with citations.
\end{enumerate}

Given an input query \(q\), the pipeline produces a set of candidate events,
a set of candidate rules, and a set of causal paths. The paths are then
verified to select a primary path. The final answer is generated from the
primary path, the verification decision, and the retained legal evidence.

\begin{figure*}[t]
    \centering
    % Replace with the official pipeline figure when the paper is finalized:
    % \includegraphics[width=0.95\textwidth]{figures/causalrag_pipeline.pdf}
    \fbox{
        \parbox{0.92\textwidth}{
        \centering
        Legal text
        \(\rightarrow\) Rule normalization
        \(\rightarrow\) Legal causal graph and causal memory
        \(\rightarrow\) Multi-hop retrieval
        \(\rightarrow\) Query-aware verification
        \(\rightarrow\) Answer and citation generation
        }
    }
    \caption{Overall architecture of the CausalRAG pipeline for reasoning
    over Vietnamese criminal law.}
    \label{fig:pipeline}
\end{figure*}

\subsection{Legal Rule Normalization}

The initial input consists of statutory provisions in textual form. The goal
of the normalization stage is to transform the text into units that can be
used for retrieval and graph reasoning. Each rule is stored with the
following principal fields:

\begin{itemize}
    \item rule identifier;
    \item statutory provision identifier and title;
    \item legal subject;
    \item applicability condition;
    \item legal effect;
    \item event ID and event name for the condition;
    \item event ID and event name for the effect;
    \item causal-relation type or effect type, when available.
\end{itemize}

An event ID is normalized as an unaccented string whose tokens are joined by
underscores. The event name retains a human-readable Vietnamese
representation. Separating the ID from the name allows the system to use a
stable ID for graph matching while using the name for semantic retrieval and
answer generation.

For example, a rule may be normalized as:

\begin{quote}
Condition: ``intentionally committing an offense but failing to complete it
for reasons beyond the offender's control'';

Effect: ``is classified as an attempted crime'';

Authority: Article 15 of the Penal Code.
\end{quote}

This rule creates the relation:

\[
\texttt{co\_y\_thuc\_hien\_toi\_pham\_nhung\_khong\_hoan\_thanh}
\rightarrow
\texttt{pham\_toi\_chua\_dat}.
\]

If an event is both the effect of a preceding rule and the condition of a
subsequent rule, the two rules can be connected into a multi-hop causal path.

\subsection{Construction of the Heterogeneous Legal Causal Graph}

From the normalized rule set, the system creates three principal groups of
nodes: events, rules, and statutory provisions. Each rule is connected to
its condition event, effect event, and source provision. In addition, a
\texttt{CAUSES} edge is created from the condition event to the effect event.

Representing rules as separate nodes preserves legal information that a
simple event-to-event edge could lose, including:

\begin{itemize}
    \item multiple rules connecting the same pair of events;
    \item rules with the same outcome but different conditions;
    \item the same event appearing in multiple statutory provisions;
    \item a relation supported by multiple legal authorities.
\end{itemize}

For graph traversal, the system uses a projection over events and
\texttt{CAUSES} edges. Each edge nevertheless retains its rule ID and article
ID so that a path can be traced back to its original evidence.

The system also extracts two-hop chains from the complete graph:

\[
v_0 \rightarrow v_1 \rightarrow v_2.
\]

These chains are used to inspect graph structure, support benchmark
construction, and generate candidate paths during retrieval.

\subsection{Construction of Causal Memory}

Graph traversal requires a seed event relevant to the query. However, an
event name in the graph may not occur verbatim in the question. The system
therefore constructs a causal memory containing both rule records and event
records.

Each rule record includes:

\begin{itemize}
    \item condition and effect text;
    \item legal subject;
    \item statutory provision;
    \item condition event and effect event;
    \item metadata linking the record to the graph.
\end{itemize}

Each event record includes:

\begin{itemize}
    \item event ID and event name;
    \item the event's role as a condition, mediator, or effect;
    \item related rules and statutory provisions;
    \item the corresponding graph node ID.
\end{itemize}

The records are converted into vector representations using BGE-M3
\cite{chen2024bgem3}. A vector index is constructed with FAISS to support
nearest-neighbor search over the memory \cite{johnson2019billion}. In the
experiments, the same embedding space is used for queries, rule text, and
event text.

For a query \(q\), the query embedding is computed as:

\begin{equation}
\mathbf{z}_q
=
f_{\theta}(q),
\end{equation}

where \(f_{\theta}\) is the BGE-M3 encoder. The semantic similarity between
the query and memory record \(m_i\) is computed using cosine similarity:

\begin{equation}
s_{\mathrm{sem}}
\left(
q,m_i
\right)
=
\frac{
    \mathbf{z}_q^\top \mathbf{z}_{m_i}
}{
    \left\Vert\mathbf{z}_q\right\Vert
    \left\Vert\mathbf{z}_{m_i}\right\Vert
}.
\end{equation}

The semantic score is used only to generate and rank candidates. A record
with a high semantic score does not automatically become final evidence.

\subsection{Candidate Event and Rule Retrieval}

The retrieval stage operates through two parallel branches.

The first branch retrieves events with high similarity to the query. These
events are used as seeds for graph traversal. Each seed record stores:

\begin{itemize}
    \item event ID;
    \item graph node ID;
    \item event name;
    \item semantic score;
    \item event role;
    \item retrieval method.
\end{itemize}

The second branch directly retrieves rules whose content is similar to the
query. Direct rule retrieval retains relevant rules even when their event
names do not closely match the phrasing of the question.

Candidates with scores below the minimum threshold are discarded. In the
main configuration, the system selects at most eight seed events and eight
direct rule hits before graph expansion. A larger semantic pool is retained
temporarily to avoid prematurely discarding candidates that may participate
in a strong path.

\subsection{Multi-Hop Causal-Path Expansion}

From each seed event, the retriever traverses \texttt{CAUSES} edges in the
configured direction. The experiments use \texttt{direction=both}, enabling
both forward and reverse search. Depending on the position of an event in
the query, a seed may serve as a cause, mediator, or outcome.

Expansion is constrained by:

\begin{itemize}
    \item the maximum number of hops;
    \item the maximum number of paths per seed event;
    \item the maximum number of candidate rules;
    \item the semantic-score thresholds for events and rules.
\end{itemize}

Each candidate path stores:

\begin{itemize}
    \item the ordered lists of event IDs and event names;
    \item the list of rule IDs supporting the individual hops;
    \item the list of article IDs;
    \item the seed event and outcome event;
    \item the number of hops;
    \item semantic and graph scores;
    \item a structural explanation of the path.
\end{itemize}

The retriever prioritizes paths that are structurally continuous, have a
complete rule for every hop, and are highly relevant to the query. Duplicate
paths are merged according to their event and rule sequences. The final
retrieval result contains candidate paths, rule evidence, and retrieved
events.

This design separates two concepts:

\begin{itemize}
    \item \emph{path retrieval}: whether the gold path appears in the
    candidate set;

    \item \emph{path ranking}: whether the gold path is ranked first.
\end{itemize}

This distinction is reflected by the Oracle Exact Path and Top-1 Exact Path
metrics in Section~\ref{sec:results}.

\subsection{Query Analysis by Claim Type}

The retrieval result is passed to query-aware verification. Before paths are
verified, the system normalizes the query by converting it to lowercase,
removing redundant whitespace, and recognizing groups of linguistic signals.

The main signal groups include:

\begin{itemize}
    \item signals of a direct relation, such as ``directly'' or ``direct
    edge'';

    \item signals indicating mediator removal or an assumption that the
    mediator is absent;

    \item signals that the outcome remains;

    \item signals that the outcome disappears;

    \item signals asking about necessity;

    \item counterfactual or hypothetical-condition signals.
\end{itemize}

From these signals, the query analyzer creates a
\texttt{query\_analysis} object containing:

\begin{itemize}
    \item the normalized query;
    \item the claim type;
    \item the intent-detector version;
    \item the direct, remove, remains, disappears, and necessary flags;
    \item the mapped mediator;
    \item the factual outcome and counterfactual outcome, when available;
    \item the verification method;
    \item the structural-fallback status.
\end{itemize}

For counterfactual questions, the mediator is mapped to an event on a
candidate path using exact matching or semantic matching. If no mediator can
be identified with sufficient confidence, the system does not automatically
assume an arbitrary event; instead, it retains an undetermined result or
uses a fallback while recording the reason.

\subsection{LegalSCM Initialization}

LegalSCM is constructed directly from the normalized rule set. Each event is
mapped to a three-valued endogenous variable:

\[
X_v
\in
\left\{
    \mathrm{TRUE},
    \mathrm{FALSE},
    \mathrm{UNKNOWN}
\right\}.
\]

Each rule is converted into a mechanism linking a condition event to an
effect event. In the current corpus version, a basic mechanism has the form:

\begin{equation}
X_{e_i}
\leftarrow
f_{r_i}
\left(
    X_{c_i}
\right).
\end{equation}

When \(X_{c_i}=\mathrm{TRUE}\), rule \(r_i\) can activate and infer the effect
event. If the condition is undetermined, the rule lacks sufficient support
to activate. LegalSCM iterates over the mechanisms until reaching a fixed
point or the maximum number of iterations.

Root events not assigned by the factual context remain in the
\(\mathrm{UNKNOWN}\) state. For endogenous events, the system applies a
controlled closed-world assumption when no mechanism can activate them. If
mechanisms produce conflicting results, the final state is set to
\(\mathrm{UNKNOWN}\) rather than forcing a choice.

LegalSCM is cached according to the rule file and its modification time. The
cache prevents the complete set of 2,884 rules from being parsed again for
every question in a batch.

\subsection{Factual Causal-Path Validation}

For a candidate path:

\[
p=
\left(
v_0,r_1,v_1,\ldots,r_h,v_h
\right),
\]

the system places seed event \(v_0\) in the factual context and performs
fixed-point inference. A path is considered factually supported only when:

\begin{enumerate}
    \item outcome event \(v_h\) is inferred to be in the
    \(\mathrm{TRUE}\) state;

    \item the rules expected along the path occur in the set of activated
    rules;

    \item each hop has at least one matching mechanism;

    \item no conflict causes a critical event to become
    \(\mathrm{UNKNOWN}\).
\end{enumerate}

The rule coverage of a path is computed as:

\begin{equation}
\mathrm{Coverage}(p)
=
\frac{
    \left|
        \mathcal{R}_{p}
        \cap
        \mathcal{R}_{\mathrm{activated}}
    \right|
}{
    \left|
        \mathcal{R}_{p}
    \right|
},
\end{equation}

where \(\mathcal{R}_{p}\) is the set of rules expected along the path and
\(\mathcal{R}_{\mathrm{activated}}\) is the set of rules activated in the
factual world.

In addition to aggregate coverage, the system records the coverage state of
each hop. This trace distinguishes cases in which the outcome is inferred
through the path under consideration from cases in which the outcome is
activated by another rule outside that path.

\subsection{Hard Intervention on a Mediator}

For each mediator \(M\) within a multi-hop path, LegalSCM constructs a
counterfactual world through the operation:

\[
do
\left(
    M=\mathrm{FALSE}
\right).
\]

In this world, the ordinary mechanism that generates \(M\) is replaced by
the constant assignment \(\mathrm{FALSE}\). Rules whose conditions depend on
\(M\) are reevaluated, after which the states of descendant events are
re-inferred to a fixed point. This implementation follows the intervention
semantics of SCMs \cite{pearl2009causality}.

For outcome \(Y\), the system compares:

\[
Y^{F}
=
Y
\left(
W^{F}
\right)
\]

with:

\[
Y^{CF}
=
Y
\left(
W^{F}
\mid
do(M=\mathrm{FALSE})
\right).
\]

If \(Y^{F}=\mathrm{TRUE}\) and \(Y^{CF}=\mathrm{FALSE}\), the mediator is
assigned the \texttt{NECESSARY} status in the context under consideration.
If the outcome remains \(\mathrm{TRUE}\), the mediator is not wholly
necessary within the model. If the factual or counterfactual outcome is
undetermined, the intervention is assigned an unresolved status.

Each intervention trace records:

\begin{itemize}
    \item mediator ID and mediator name;
    \item intervention assignment;
    \item factual mediator state;
    \item factual outcome;
    \item counterfactual outcome;
    \item whether the outcome changes;
    \item disabled rules;
    \item newly activated rules;
    \item alternative rules leading to the same outcome;
    \item intervention status and explanation.
\end{itemize}

These fields are stored in the Step 4 output for auditing and are not used
merely to compute a single decision label.

\subsection{Node-Deletion Path Ablation}

Path ablation is implemented by removing the mediator from the event graph
and then searching for alternative paths from the seed event to the outcome.
The search for alternative paths is constrained by the maximum number of
hops and paths to control cost.

For each alternative path, the system records:

\begin{itemize}
    \item the alternative event sequence;
    \item the number of hops;
    \item supporting rules;
    \item graph score;
    \item the score of the best alternative path.
\end{itemize}

If no alternative path is found, the mediator is considered necessary under
the reachability criterion. If an alternative path exceeds the score
threshold, the outcome is considered still reachable after the mediator is
deleted.

When the system runs in \texttt{structural\_scm} mode, the path-ablation
result is still computed and stored as a baseline diagnostic. If LegalSCM
cannot be initialized or cannot verify a path, the system may use path
ablation as a fallback, but it must explicitly record:

\begin{itemize}
    \item \texttt{structural\_fallback\_used};
    \item the reason for the fallback;
    \item the method that actually produced the result.
\end{itemize}

In the paired experiment, \texttt{path\_ablation} mode is run independently
with the same retrieval output and hyperparameters, enabling the two engines
to be compared without changing any other component.

\subsection{Evidence Ranking and Classification}

Each rule evidence item receives a verification score that combines four
signals:

\begin{itemize}
    \item the degree of support from factual paths;
    \item the degree of support from counterfactual verification;
    \item semantic relevance to the query;
    \item graph relevance and retrieval score.
\end{itemize}

Evidence on the primary causal path receives an additional priority score.
Based on the verification score, evidence is divided into three groups:

\begin{description}
    \item[\texttt{KEEP}] Evidence strong enough to be passed to the answer
    generation stage.

    \item[\texttt{UNCERTAIN}] Evidence that is relevant but not sufficiently
    supported.

    \item[\texttt{REMOVE}] Evidence that does not meet the threshold or is
    rejected by the verification result.
\end{description}

In the experimental configuration, the keep threshold is set to \(0.52\),
whereas the reject threshold is \(0.34\). Evidence on the primary path is
retained in the path's rule order rather than being sorted only by
verification score.

\subsection{Primary-Path Selection and Final Decision}

Paths are ranked according to:

\begin{itemize}
    \item factual-support status;
    \item path completeness;
    \item number of hops;
    \item consistency score;
    \item graph and retrieval scores;
    \item compatibility with the claim type.
\end{itemize}

The system prioritizes paths with \texttt{SUPPORTED} status, complete two-hop
structure for multi-hop questions, and high consistency scores. The
highest-ranked path is selected as the primary path, while the remaining
paths are retained as supporting paths or for branch/convergence reasoning.

Unlike an earlier majority-over-paths decision rule, the query-aware version
produces the final decision from the claim type and the primary path.
Specifically:

\begin{itemize}
    \item for a standard causal claim, the system checks the factual support
    of the primary path;

    \item for a direct claim, the system checks for a direct edge rather than
    merely checking reachability;

    \item for a mediator claim, the system compares factual and
    counterfactual outcomes;

    \item if the mediator mapping or a world state is undetermined, the
    decision becomes \texttt{UNCERTAIN}.
\end{itemize}

The verification output includes:

\begin{itemize}
    \item \texttt{final\_decision};
    \item \texttt{decision\_score};
    \item \texttt{decision\_explanation};
    \item \texttt{primary\_path\_ids};
    \item \texttt{query\_analysis};
    \item \texttt{path\_verifications};
    \item retained, uncertain, and removed evidence;
    \item the Step 4 configuration and version.
\end{itemize}

\subsection{Answer-Intent Inference}

The answer-generation stage receives the complete verification output rather
than only a list of rules. The answer intent is inferred from the query text
and \texttt{query\_analysis}. Direct or counterfactual claims are always
preferentially mapped to \texttt{CLAIM\_VERIFICATION} to preserve the
semantics of Step 4.

For the remaining questions, the system recognizes linguistic patterns
related to:

\begin{itemize}
    \item the final effect;
    \item the initial cause;
    \item the bridge event;
    \item chain explanation;
    \item yes/no questions;
    \item branching;
    \item convergence.
\end{itemize}

This process does not use \texttt{question\_type}. Therefore, the same
mechanism can be applied to queries outside the benchmark when their phrasing
contains the corresponding signals.

\subsection{Path and Evidence Selection for the Answer}

Step 5 selects paths first and then selects evidence according to the chosen
paths. The primary path is always retained when one exists.

For linear questions, evidence on the primary path is placed first and
retained in causal order. For branch reasoning, the system retains paths
that share a common prefix with the primary path. For convergence reasoning,
the system retains paths with a common suffix. This mechanism avoids the
loss of branches required for the answer that would result from selecting
only one path.

The final evidence is selected from:

\begin{enumerate}
    \item rules on the primary path;
    \item rules on structurally related paths;
    \item supplementary evidence verified by Step 4;
    \item retrieval-evidence fallback if verification retains no evidence.
\end{enumerate}

Duplicate rules are removed by rule ID. When citations are created, evidence
items representing the same causal edge are further deduplicated to avoid
repeatedly listing equivalent authorities.

\subsection{Extractive Answer Generation}

Although the implementation supports external providers such as Ollama,
OpenAI-compatible APIs, and Gemini, the main experiments use the extractive
provider. This choice makes the results deterministic and separates the
effects of retrieval and verification from the variability of a generative
model.

The extractive provider generates an answer according to the answer intent:

\begin{description}
    \item[\texttt{FINAL\_OUTCOME}] Use the complete \texttt{effect} content
    of the final rule on the primary path. If no corresponding rule can be
    found, the system falls back to the outcome event name.

    \item[\texttt{INITIAL\_CAUSE}] Return the first event and the complete
    causal chain.

    \item[\texttt{BRIDGE\_EVENT}] Return one or more mediators between the
    seed and the outcome.

    \item[\texttt{CHAIN\_EXPLANATION}] Present each hop in order, with a
    citation to the rule supporting that hop.

    \item[\texttt{YES\_NO}] Begin with ``Yes,'' ``No,'' or ``Insufficient
    evidence,'' consistently with the final decision.

    \item[\texttt{BRANCH}] Aggregate the outcomes of paths with a common
    prefix.

    \item[\texttt{CONVERGENCE}] Aggregate the input conditions of paths with
    a common suffix.

    \item[\texttt{CLAIM\_VERIFICATION}] Answer the claim directly and use
    the decision explanation from Step 4.
\end{description}

For example, consider the primary path:

\[
\text{Intentional commission of an offense without completion}
\rightarrow
\text{Attempted crime}
\rightarrow
\text{Criminal liability},
\]

where the final rule contains the effect ``must bear criminal liability for
the attempted offense.'' The answer therefore uses this complete effect
rather than only the general event label ``Criminal liability.''

\subsection{Citation and Output Validation}

Each evidence item is assigned an index \([E_i]\). The answer validator
extracts citations appearing in the answer and checks them against the list
of selected evidence.

The validation procedure includes:

\begin{itemize}
    \item removing nonexistent citations;
    \item prioritizing citations on the primary path;
    \item adding a citation when the answer uses evidence but lacks a marker;
    \item recording validation warnings;
    \item mapping evidence citations to article IDs for evaluation.
\end{itemize}

The final output stores both evidence-index citations and authorities of the
form ``Article \(X\).'' This enables checks of both the answer's internal
grounding and its accuracy against gold citations.

\subsection{Provenance and Reproducibility}

Each artifact records the version of the stage that produced it. In the main
experiment:

\begin{itemize}
    \item Step 4 uses version
    \texttt{4.1-query-aware-legal-scm};

    \item Step 5 uses version
    \texttt{5.1-query-aware-grounded-answer};

    \item the batch runner uses version
    \texttt{2.1-query-aware-answer};

    \item the evaluator uses version
    \texttt{2.0-current-pipeline-schema}.
\end{itemize}

The generation metadata also records:

\begin{itemize}
    \item answer intent;
    \item intent-detector version;
    \item counterfactual-signal source;
    \item Step 4 and Step 5 versions;
    \item structural result or structural fallback;
    \item provider and model;
    \item execution time;
    \item the number of evidence items and paths used;
    \item the list of rules on the primary path.
\end{itemize}

This provenance allows results from different pipeline versions to be
distinguished, preventing old and new artifacts from being compared under
the same experimental label.
